% =============================================================
% §4 System Design
% =============================================================
\section{System Design}
\label{sec:system}

Our system has two co-designed components: an AIGC inference
simulator that models the five physics from \S\ref{sec:intro}
(\S\ref{sec:simulator}), and an AIGC-aware PPO scheduler that
exposes those physics to the policy via structured state and
reward shaping (\S\ref{sec:rl}).

\subsection{AIGC Inference Simulator}
\label{sec:simulator}

We implement the simulator in Python as a discrete-time
event-driven environment with a 0.1~s step. Each step
(a)~releases completed tasks, (b)~promotes
\texttt{WAITING}~$\to$~\texttt{READY} as dependencies and
arrival times allow, (c)~invokes the scheduler, and
(d)~starts admitted tasks and accumulates per-server energy.
The five AIGC physics are realized as four engineering
milestones M1--M4.

\textbf{M1: Model weight residency.} Each server tracks
\texttt{loaded\_models} with reference counts; an unloaded
model triggers LRU eviction of unpinned models and a
\texttt{cold\_load\_sec} delay (5~s for LLaMA-7B / 14~GB,
25~s for LLaMA-70B-INT8 / 70~GB; calibrated to NVMe
sequential read).

\textbf{M2: Prefill--Decode two-phase tasks with KV cache.}
Each request decomposes into (Prefill, Decode) sharing
\texttt{req\_id}; per-phase workload is
\texttt{tokens$\cdot$prefill\_tflops}/$k$ and the analogous
decode form. We encode KV size into
\texttt{prefill.output\_size}
($\text{prompt\_tokens}\cdot\kappa_m/1024$~GB), so the existing
transfer-time formula automatically charges
\textbf{KV-migration cost} whenever Decode lands on a server
other than where Prefill ran---no extra code paths. Each task
also carries \texttt{kv\_cache\_GB} occupying VRAM during
execution.

\textbf{M3: Continuous batching.} An admitted task counts the
running same-model same-kind tasks
(\texttt{batch\_size\_at\_admit}) and pays effective time
$T_{\text{solo}}(1+(\text{batch\_size}-1)\cdot\beta_k)$ with
$\beta_{\text{decode}}{=}0.05$ and
$\beta_{\text{prefill}}{=}0.30$ (matched to vLLM). The
\texttt{--ablation no\_batching} switch disables batching
\emph{and} serializes inference on each GPU---modeling an
engine without iteration-level batching.

\textbf{M4: Memory-bandwidth floor + realistic workloads.}
Decode is memory-bandwidth-bound; we cap effective time by a
per-model token floor $\phi_m$ (20~ms/tok decode for
LLaMA-7B, scaling to 80~ms/tok for LLaMA-70B). M4 also adds
Poisson arrival (\texttt{arrival\_time} gates
\texttt{READY}) and log-normal prompt/output lengths
calibrated to Azure LLM Inference Trace 2023
(prompt $\mu_{\log}{=}5.5,\sigma{=}1.0$;
output $\mu_{\log}{=}4.5,\sigma{=}1.2$).

\textbf{Energy model.} Each server's instantaneous power is
$P(t) = \text{idle\_W} + (\text{max\_W}-\text{idle\_W})\cdot
u_s(t)$ with per-tier coefficients calibrated to NVIDIA TDPs
(cloud A100: 50--400~W; edge A100: 30--250~W; T4: 20--70~W;
Jetson AGX Orin: 10--30~W). System energy per token is the
integrated power divided by completed output tokens.

The simulator exposes a one-call
\texttt{scheduler.schedule()} interface (read-only access to
tasks/servers/network; dispatch via
\texttt{server.add\_task}), so any scheduler---from
50-line RoundRobin to our PPO agent---can be swapped in
without environment changes. Fig.~\ref{fig:simulator-arch}
shows the end-to-end data flow.

\begin{figure*}[t]
  \centering
  \includegraphics[width=0.92\textwidth]{trace_integration.pdf}
  \caption{End-to-end architecture of our AIGC inference simulator.
    Color bands denote (purple) trace sources, (blue) workload
    generator, (red) simulation environment with M1--M4 physics
    modules, (yellow) metrics collection, and (gray) output
    artifacts. The minimal \texttt{scheduler.schedule()} interface
    decouples policy from environment, letting all nine baselines
    and twelve ablation configurations share the same physics
    backbone for fair comparison.}
  \label{fig:simulator-arch}
\end{figure*}

\subsection{AIGC-Aware Reinforcement Learning Scheduler}
\label{sec:rl}

We cast per-task dispatch as an MDP and solve it with PPO;
Fig.~\ref{fig:rl-arch} overviews the resulting decision loop,
from per-task state encoding through the actor-critic policy to
the seven-term reward and the PPO update that closes the loop.

\begin{figure*}[t]
  \centering
  \includegraphics[width=0.86\textwidth]{rl_scheduler_arch.pdf}
  \caption{AIGC-aware PPO scheduler as an MDP decision loop. For
    each \texttt{READY} task, the state encoder (\textcircled{1})
    builds a 96-dim vector whose per-server block carries five
    AIGC-specific features (highlighted); the actor-critic policy
    (\textcircled{2}) emits a masked action distribution and a
    value estimate; the chosen server is dispatched into the
    simulator (\textcircled{3}, the five AIGC physics); the
    seven-term dense reward (\textcircled{4}) groups
    load-balancing, AIGC-aware, and anti-collapse signals; and the
    PPO update (\textcircled{5}) closes the loop via policy
    gradient. Pretraining is what lets the policy escape the
    always-cloud attractor.}
  \label{fig:rl-arch}
\end{figure*}

\textbf{State.} For each \texttt{READY} task the encoder emits
a flat vector of dimension $(10+|\mathcal{M}|)+2+10N$ holding
three blocks: a \emph{task block} (compute demand, workload,
input/output size, dependency/successor counts, task-kind
one-hot, normalized KV-cache size, and a model-ID one-hot),
a 2-dim \emph{global block} (\texttt{READY} queue size,
completed ratio), and a $10$-dim per-server block of five
general features (compute/memory utilization, bandwidth,
queue length, transfer-time estimate) followed by five
\textbf{AIGC-specific features}: (i)~loaded-model
indicator, (ii)~normalized batch occupancy, (iii)~free VRAM
fraction, (iv)~normalized cold-load cost, and
(v)~\emph{sibling-server indicator}---for Decode, whether
the candidate server is where Prefill ran. The sibling
indicator is novel: it gives the policy direct read-access
to KV-cache locality without reasoning about request IDs.
For edge${=}5$, $|\mathcal{M}|{=}4$, the state dimension is
96.

\textbf{Action.} Discrete choice over the $N$ servers with
\textbf{action masking} (probability zero whenever
\texttt{can\_allocate} fails); the mask is recomputed at
every dispatch.

\textbf{Reward.} A weighted sum of seven dense per-decision
components,
\begin{equation}
  \label{eq:reward}
  \begin{aligned}
    r =\;& 0.25\,r_{\text{time}}
        + 0.10\,r_{\text{balance}}
        + 0.10\,r_{\text{match}} \\
       &+ 0.15\,r_{\text{warm}}
        + 0.20\,r_{\text{batch}}
        + 0.10\,r_{\text{affinity}}
        + 0.10\,r_{\text{cloud}},
  \end{aligned}
\end{equation}
combining three load-balancing signals
($r_{\text{time}}, r_{\text{balance}}, r_{\text{match}}$
penalize long completion, server overload, and capacity
mismatch) with four AIGC-specific signals: $r_{\text{warm}}$
($+1$ if the model is already loaded, else
$-\text{cold\_load}/20$ clipped to $[-1,0]$);
$r_{\text{batch}}$ rewards joining existing same-model
batches (weighted higher for Decode);
$r_{\text{affinity}}$ rewards co-locating Decode with its
Prefill ($+1$) and penalizes non-trivial KV migration
($-0.5$); and $r_{\text{cloud}}$ ($-u_{\text{cloud}}$
when an edge alternative exists) discourages an
``always-cloud'' attractor we observed in early experiments
(see \S\ref{sec:implications}).

\textbf{Architecture and training.} A 3-layer MLP encoder
($\text{state} \to 256 \to 128 \to 64$, LayerNorm + ReLU)
feeds an Actor (Linear $\to$ softmax-with-mask) and a Critic
(Linear $\to$ ReLU $\to$ Linear). We train with
PPO~\cite{schulman2017ppo} (clip $\epsilon{=}0.2$),
GAE~\cite{schulman2015gae} ($\gamma{=}0.95$,
$\lambda{=}0.90$), and entropy regularization ($0.05$,
reduced to $0.01$ after pretraining); each update sees 32
buffered transitions across 4 epochs. Crucially, we precede
online training with \textbf{20 episodes of pretrain}
(${\approx}$10--15~s wall-clock) initialized by 30 uniform
random warm-up steps; without this, the policy collapses to
``always-cloud'' and \slo{} drops $38\%$ (\S\ref{sec:ablation}).
Each AIGC component is gated by a constructor flag, enabling
the systematic 12-component ablation in
\S\ref{sec:ablation}.
